\documentclass{article}
\usepackage[a4paper, margin=1in]{geometry} % Adjust margins
\usepackage{amsmath, amssymb, amsthm}
\usepackage{tikz}
\usepackage{xcolor}
\usepackage{graphicx}
\usepackage{listings}
\usepackage{fancyhdr}
\usetikzlibrary{automata, positioning}
\usepackage{helvet} % Match font to reference
\pagestyle{fancy}
\usepackage{enumitem,xcolor}


% Ensure proper header spacing
\setlength{\headheight}{15pt} % Adjust height of the header
\setlength{\headsep}{10pt} 

% Define Header/Footer
\fancyhf{}
\fancyhead[L]{\textbf{Advanced Topics in Computer Science}} % Left header
\fancyhead[R]{DFA Minimization} % Right header
\cfoot{\thepage} % Page number at center bottom

% Define colors for section and subsection titles
\usepackage{titlesec}
\definecolor{sectioncolor}{rgb}{0.1, 0.3, 0.7} % Dark blue
\titleformat{\section}{\large\bfseries\color{sectioncolor}}{\thesection}{1em}{}
\titleformat{\subsection}{\bfseries\color{sectioncolor}}{\thesubsection}{1em}{}

% Define custom colors for different text types
\usepackage{tcolorbox}
\tcbuselibrary{theorems}
\newcommand{\important}[1]{\textcolor{red}{#1}}
\newtcbtheorem
  []% init options
  {definition}% name
  {Definition}% title
  {%
    colback=blue!5,
    colframe=blue!35!black,
    fonttitle=\bfseries,
  }% options
  {def}% prefix
\newcommand{\note}[1]{\textcolor{green!50!black}{#1}}

% Theorem environment setup
\newtheorem{theorem}{Theorem}
\newtheorem{lemma}{Lemma}
\newtheorem{corollary}{Corollary}

\definecolor{codegray}{gray}{0.9}
\lstdefinestyle{customc}{
  backgroundcolor=\color{codegray},
  language=C,
  basicstyle=\ttfamily\footnotesize,
  keywordstyle=\color{blue},
  commentstyle=\color{green!50!black},
  stringstyle=\color{red}
}
\begin{document}

\title{\textbf{DFA Minimizer}}
\author{Doodlebob}
\date{}
\maketitle
\thispagestyle{fancy}

The following is an auto-explainer for DFA Minimization as represented by a series of figures depicting the evolving partitions of the DFA’s state set. \\

A Deterministic Finite Automaton (DFA) is denoted by the quintuple
\[
M = (Q, \Sigma, \delta, q_0, F)
\]
and the minimization process works by iteratively refining partitions of the state set until all equivalent states are grouped together. In the minimization process:
\begin{itemize}
    \item The initial partition divides $Q$ into two blocks: the final states $F$ and the non-final states $Q \setminus F$.
    \item At each subsequent step, the partition is refined by examining the transitions for every input symbol and grouping together states that behave equivalently.
    \item This iterative refinement continues until no further splitting is possible, thus yielding a minimal DFA.
\end{itemize}

\section{DFA Definition and Minimization Overview}

A Deterministic Finite Automaton is defined as a 5-tuple 
\[
(Q, \Sigma, \delta, q_0, F)
\]
where:
\begin{itemize}
    \item $Q$ is a finite set of states,
    \item $\Sigma$ is a finite input alphabet,
    \item $\delta : Q \times \Sigma \rightarrow Q$ is the transition function,
    \item $q_0 \in Q$ is the start state, and
    \item $F \subseteq Q$ is the set of final (accepting) states.
\end{itemize}

The minimization of a DFA entails computing an equivalent DFA with the fewest possible states. The process begins by partitioning the state set into final and non-final states and then repeatedly refining these partitions by identifying states that are indistinguishable with respect to the input symbols. The resulting minimized DFA retains the same language as the original.
\section{DFA Transition Function}

The DFA's transition function, $\delta$, is defined as follows:
\begin{align*}\delta(1,\,a) &= 2\\[0.5em]\delta(1,\,b) &= 4\\[1em]\delta(2,\,a) &= 4\\[0.5em]\delta(2,\,b) &= 3\\[1em]\delta(3,\,a) &= 3\\[0.5em]\delta(3,\,b) &= 3\\[1em]\delta(4,\,a) &= 2\\[0.5em]\delta(4,\,b) &= 5\\[1em]\delta(5,\,a) &= 6\\[0.5em]\delta(5,\,b) &= 5\\[1em]\delta(6,\,a) &= 4\\[0.5em]\delta(6,\,b) &= Z\\[1em]\delta(Z,\,a) &= Z\\[0.5em]\delta(Z,\,b) &= Z
\end{align*}

\section{DFA Minimization Visualizations}

The following sequence of minimization steps was performed on the DFA. In each step, the current partition of the DFA's states is visualized. States grouped in the same block are considered equivalent with respect to their transitions. The process starts with the initial partition (separating final and non-final states) and iteratively refines until no further changes occur.
\subsection*{Minimization Step 1}
\textbf{Explanation:}
\newline Initially, the DFA's states are partitioned into final and non-final states. The partition is as follows:
\begin{itemize}
    \item \textbf{Final States:} \textbf{ ['3'] }
    \item \textbf{Non-Final States:} \textbf{ ['1', '2', '4', '5', '6', 'Z'] }
\end{itemize}

\begin{figure}[h]
    \centering
    \includegraphics[width=0.8\linewidth]{plots/plot1.png}
    \caption{Partition configuration at minimization step 1}
\end{figure}
\subsection*{Minimization Step 2}
\textbf{Explanation:}
\newline In this minimization step, the partition is refined based on the transitions for each input symbol.
\newline \textbf{Current Partition:} \textbf{ [['3'], ['1', '4', '5', '6', 'Z'], ['2']] }\\[1ex]
\textbf{Description:} Step 2: Partitioned into 3 block(s) -- 3; 1, 4, 5, 6, Z; 2

\begin{figure}[h]
    \centering
    \includegraphics[width=0.8\linewidth]{plots/plot2.png}
    \caption{Partition configuration at minimization step 2}
\end{figure}
\subsection*{Minimization Step 3}
\textbf{Explanation:}
\newline In this minimization step, the partition is refined based on the transitions for each input symbol.
\newline \textbf{Current Partition:} \textbf{ [['3'], ['1', '4'], ['5', '6', 'Z'], ['2']] }\\[1ex]
\textbf{Description:} Step 3: Partitioned into 4 block(s) -- 3; 1, 4; 5, 6, Z; 2

\begin{figure}[h]
    \centering
    \includegraphics[width=0.8\linewidth]{plots/plot3.png}
    \caption{Partition configuration at minimization step 3}
\end{figure}
\subsection*{Minimization Step 4}
\textbf{Explanation:}
\newline In this minimization step, the partition is refined based on the transitions for each input symbol.
\newline \textbf{Current Partition:} \textbf{ [['3'], ['1'], ['4'], ['5', 'Z'], ['6'], ['2']] }\\[1ex]
\textbf{Description:} Step 4: Partitioned into 6 block(s) -- 3; 1; 4; 5, Z; 6; 2

\begin{figure}[h]
    \centering
    \includegraphics[width=0.8\linewidth]{plots/plot4.png}
    \caption{Partition configuration at minimization step 4}
\end{figure}
\subsection*{Minimization Step 5}
\textbf{Explanation:}
\newline In this minimization step, the partition is refined based on the transitions for each input symbol.
\newline \textbf{Current Partition:} \textbf{ [['3'], ['1'], ['4'], ['Z'], ['5'], ['6'], ['2']] }\\[1ex]
\textbf{Description:} Step 5: Partitioned into 7 block(s) -- 3; 1; 4; Z; 5; 6; 2

\begin{figure}[h]
    \centering
    \includegraphics[width=0.8\linewidth]{plots/plot5.png}
    \caption{Partition configuration at minimization step 5}
\end{figure}
\end{document}